\section{Matrices de moments et réalisations amplituédriques de rang supérieur}
\label{sec:higher-rank}
La partition du registre dodécaphasique permet de prolonger la construction de rang un à des points explicites de rang supérieur. Les mêmes treize nœuds ordonnés fournissent simultanément des mineurs de Plücker positifs et une forme de moments définie positive. Ces deux positivités sont obtenues à partir d'une seule liste d'intervalles ; la seconde prouve que la composition conserve le rang requis.

Soit \(d_j=K_j/Q\), où \(Q=\sum_{j=1}^{12}K_j=6263\), et définissons
\[
 t_0=0,\qquad t_j=\sum_{r=1}^jd_r.
\]
La positivité des douze blocs donne
\(0=t_0<t_1<\cdots<t_{12}=1\). Pour chaque \(2\le k\le9\), nous posons
\[
 C^{(k)}_{ai}=t_i^a\quad(0\le a<k),\qquad
 Z^{(k)}_{bi}=t_i^b\quad(0\le b<k+4),\qquad 0\le i\le12.
\]
La convention \(t_i^0=1\) inclut le nœud zéro. Ainsi, \(C^{(k)}\)
est une matrice de taille \(k\times13\), tandis que \(Z^{(k)}\) a pour
taille \((k+4)\times13\).

\begin{theorem}[Positivité et rang de la réalisation des moments]
La matrice \(C^{(k)}\) représente un point de \(\Gr_{>0}(k,13)\), et tous les mineurs maximaux ordonnés de \(Z^{(k)}\) sont positifs. La matrice
\[
 Y^{(k)}=C^{(k)}(Z^{(k)})^{\mathsf T},\qquad
 Y^{(k)}_{ab}=\sum_{i=0}^{12}t_i^{a+b},
\]
est de rang \(k\). En particulier, pour la donnée externe positive
\(Z^{(k)}\) déjà construite à partir du registre,
\[
 [Y^{(k)}]\in\mathcal A_{13,k,4}(Z^{(k)}),\qquad 2\le k\le9,
\]
où
\[
 \mathcal A_{13,k,4}(Z^{(k)})
 =\{[C(Z^{(k)})^{\mathsf T}]:[C]\in\Gr_{\ge0}(k,13)\}.
\]
\end{theorem}
\begin{proof}
Un mineur maximal de l’une ou l’autre des deux matrices initiales a la forme
\[
 \det(t_{i_j}^{a})_{a=0}^{r-1}
 =\prod_{p<q}(t_{i_q}-t_{i_p})>0,
\]
pour des indices croissants. L'inégalité prouve les deux affirmations sur Plücker. Le bloc initial \(k\times k\) de \(Y^{(k)}\) est la matrice de Hankel \(H^{(k)}\). Pour \(x\ne0\),
\[
 x^{\mathsf T}H^{(k)}x
 =\sum_{i=0}^{12}\left(\sum_{a=0}^{k-1}x_at_i^a\right)^2>0.
\]
En effet, un polynôme non nul de degré inférieur à \(k\le9\) ne
s’annule pas en treize points distincts. Donc \(H^{(k)}\) est définie positive,
son déterminant est strictement positif et \(Y^{(k)}\) est de rang
\(k\). L’appartenance amplituédrique est alors la composition des
deux matrices positives précédentes, dont le rang est déjà vérifié.
\end{proof}

En rang deux, la preuve prend la forme particulièrement explicite
\[
 \det H^{(2)}=13\sum_i t_i^2-\left(\sum_i t_i\right)^2
 =\sum_{i<j}(t_j-t_i)^2>0.
\]
La séparation des nœuds, produite par la positivité du registre, est ainsi convertie en coercivité d'une forme quadratique. À chaque registre fixé correspond un point déterminé pour chaque rang. La couverture complète, les triangulations et les résidus démontrés dans la section précédente conservent leur domaine de rang un ; le présent théorème ajoute la collection explicite de points de rangs deux à neuf et leurs matrices de moments. La distinction permet de réutiliser chaque démonstration avec sa portée exacte.
